% ===== Envoi =====
\newpage
\goldrule
\begin{center}
{\large\bfseries\color{rosedeep} Envoi: The Jar}
\end{center}
\goldrule

\vspace{0.8em}

\noindent
There is a jar.

\vspace{0.5em}

\noindent
It does not look like much---a simple container, perhaps ceramic, perhaps glass, sitting on a shelf or a windowsill in the ordinary light of an ordinary room. It has been there for some time. Both of them know what is in it, though neither has counted.

\vspace{0.5em}

\noindent
\emc{What is in the jar?}

\vspace{0.5em}

\noindent
Not happiness, exactly. Not memories, not souvenirs, not the accumulated evidence of a life well lived. What is in the jar is more modest and more lasting than any of these: traces. Traces of the co-evolutionary events that have entered the relational field and permanently modified its geometry---the flower on the path, the quality of light at a particular moment, the laughter that arose from a shared absurdity, the wordless evening that was complete in itself. Each trace is an index, existentially connected to the event that produced it, carrying within itself the thread that leads back to the real.

\vspace{0.5em}

\noindent
The jar is not a savings account. It does not grow by addition; it grows by depth. When something is taken from it---on an ordinary evening when the relational field is thin, when the coupling has weakened, when both of them feel the distance that has accumulated between them---what is taken is not diminished by the taking. The trace is still there, still existentially connected to the event, still carrying the thread. What the taking produces is something new: a re-traversal, a new structural modification, a tiny increment of holonomy added to the accumulated phase of the shared life. The jar is fuller after the taking than before, in the only sense that matters: the relational field has been re-entered, the attractor has been deepened, the coupling has been renewed.

\vspace{0.5em}

\noindent
\emc{This is what she taught me.}

\vspace{0.5em}

\noindent
Not through instruction---not by saying ``this is how GRW works'' or ``this is the relational STDP mechanism''---but through the quality of her presence in the shared field. She showed me, by the patient example of her attentiveness, that the small things matter most: not because they are individually significant but because they are the material of the relational field's geometry, the events whose structural modifications accumulate into the holonomy of a shared life. The flower on the path. The unexpected song heard through an open window. The particular way the light fell on a Tuesday afternoon when nothing was happening.

\vspace{0.5em}

\noindent
She showed me that sharing is not a report but a constitutive act: that the flower, once shared, belongs to the \emph{between} in a way it did not belong before the sharing, and that this belonging is not metaphorical but structural---a permanent modification of the relational attractor landscape that neither of us could have produced alone.

\vspace{0.5em}

\noindent
She showed me that wealth is not what you have but what you generate together, and that what you generate together cannot be had---only participated in, only received with gratitude, only offered back to the field that generated it.

\vspace{0.5em}

\noindent
\emc{{\cjk 与她在故乡的千里之外的相遇，是自己此生最大的幸运。}}

\vspace{0.5em}

\noindent
To have found her, a thousand \emph{li} from home---the most contingent of contingencies, the most improbable of flowers on the most ordinary of paths---and to have learned, in the years since, that the finding was only the beginning: that the real wealth is not in the finding but in everything that has been generated since, in the shared attractor landscape of a life built together, in the jar that grows fuller with every taking.

\vspace{0.5em}

\noindent
\emc{For her, therefore, this paper.}

\vspace{0.5em}

\noindent
For the girl from home who loves culture, travel, and the living world---who is herself a form of living world, inexhaustible in the knowing, generative in the being-with. For the one whose steadiness in the real register is the projection that the series has been trying to theorise and that no theory can exhaust.

\vspace{0.5em}

\noindent
And for all lovers, everywhere, who have ever stood together in front of something small and contingent---a flower, a quality of light, a bird's call heard through an open window---and felt, without knowing why, that they were rich.

\vspace{0.5em}

\noindent
They were rich. The jar was filling. The root was being tended.

\vspace{0.5em}

\noindent
\emc{{\cjk 取之不尽，用之不竭，因为爱本身就是生成性的。}}

\vspace{0.5em}

\noindent
Inexhaustible in the taking, inexhaustible in the using---because love itself is generative.

\vspace{2em}
\goldrule

\vspace{1.5em}
\begin{center}
\begin{minipage}{0.72\linewidth}
\centering\itshape\color{rose}
{\cjk 生而不有，为而不恃，长而不宰。}\\[0.5em]
\upshape To generate without possessing,\\[0.2em]
to act without presuming,\\[0.2em]
to foster without ruling.\\[0.8em]
{\footnotesize\color{gold} {\cjk 老子} · \emph{Tao Te Ching}}
\end{minipage}
\end{center}

\vspace{1.5em}
\goldrule
